% Root, A, and D full literal audits passed. The frozen154-profile computation is preserved.
\subsection{Two equal incoming hom-set sizes at a monoid endpoint}

\begin{proposition}\label{C:two-small-incoming-homs}
Let $\mathcal C$ have objects $0,1,2$, with endomorphism monoids
$1,1,M$, respectively, and no backward arrows.  Suppose
\[
 X=\mathcal C(0,1)\ne\varnothing,
 \qquad |\mathcal C(1,2)|=|\mathcal C(0,2)|=2.
\]
Then $B\mathcal C$ is not homotopy equivalent to $\mathbb CP^2$.
No bound on the order of the finite monoid $M$ is required.
\end{proposition}

\begin{proof}
Put $T=\mathcal C(1,2)$ and $U=\mathcal C(0,2)$, regarded as left
$M$-sets.  Each $x\in X$ gives an equivariant map
$\beta_x:T\longrightarrow U$, $t\longmapsto tx$.
If one of these maps is constant, its value is fixed by every element
of $M$.  The lower full subcategory on $0,1$ is a graph, and in
particular has zero rational second cohomology.  The fixed incoming
arrow lemma therefore makes every rational degree-two product of
$B\mathcal C$ zero.

We may consequently assume that every $\beta_x$ is a bijection.
Choose $x_0\in X$ and identify $U$ with $T$ by $\beta_{x_0}$.
Every $\beta_x$ is now either the identity or the transposition
$\tau$ of this two-element set.  If both occur, $\tau$ commutes with
the $M$-action, so every actor is the identity or $\tau$.
Label the two states by $\mathbb Z/2$.  There is an actual functor
\[
 \mathcal C\longrightarrow B(\mathbb Z/2)
\]
which sends an endomorphism in $M$ to its permutation, an arrow in
$T$ or $U$ to its state label, and $x\in X$ to the parity of
$\beta_x$.  The equivariance and composition formulas verify every
functor equation.  If $\beta_x=\tau$, the loop represented by $x$
and the reverse of $x_0$ maps nontrivially.  Thus the fundamental
group of $B\mathcal C$ surjects onto $\mathbb Z/2$.

It remains that all $\beta_x$ are the same bijection.  Collapse $X$
to its selected arrow $x_0$, obtaining an actual functor
$q:\mathcal C\longrightarrow\mathcal D$ with a literal section.
Let $E$ be the action category of the left $M$-set $T$.
If $E$ is disconnected, both elements of $T$ are fixed by every
actor, and the fixed incoming arrow argument applies again.
Assume therefore that $E$ is connected.

Here the normalized rational chain kernel of $q$ can be read off
without a spectral sequence.  In degree one it has basis
$[x]-[x_0]$, for $x\ne x_0$.  In degree $n\ge2$ its basis consists
of the differences of chronological simplices
\[
 [x,t,a_1,\ldots,a_{n-2}]
 -[x_0,t,a_1,\ldots,a_{n-2}],
\]
where $t\in T$ and the $a_i$ are composable nonidentity arrows of
$M$.  The face forgetting $x$ cancels, as does the face composing
$t$ with $x$, since all $\beta_x$ coincide.  The remaining faces
are precisely the augmented action-category boundary, up to the
uniform shift sign.  Writing $\widetilde{\mathbb Q X}$ for the
augmentation kernel, this identifies the chain kernel with
\[
 \widetilde{\mathbb Q X}\otimes
 \operatorname{Cone}\bigl(C_*(BE;\mathbb Q)
                    \xrightarrow{\epsilon}\mathbb Q\bigr)[1].
\]
Since $E$ is nonempty and connected, the displayed chain kernel
has zero homology in degrees one and two.  Its long exact sequence consequently gives
an isomorphism
$q^*:H^2(B\mathcal D;\mathbb Q)\longrightarrow
H^2(B\mathcal C;\mathbb Q)$.

Finally, the single arrow $0\to1$ in $\mathcal D$ deletes object
$1$ by an actual deformation-retract functor: send $1$ to $0$,
send $t:1\to2$ to $tx_0$, and retain the monoid $M$.
The classifying space is therefore the ordinary left-act cone on
the two-element set $T$.  The universal two-cross-arrow theorem
makes all its rational degree-two products zero.  Naturality and
the preceding $H^2$ isomorphism finish the proof.
\end{proof}

\begin{proposition}\label{C:two-by-three-incoming-small-endpoint}
Let $\mathcal C$ have the same three-object support and trivial
first two corners, with $|M|\le8$, $X\ne\varnothing$,
$|\mathcal C(1,2)|=2$, and $|\mathcal C(0,2)|=3$.
Then every rational degree-two square on $B\mathcal C$ is zero.
\end{proposition}

\begin{proof}
If $H^2(BM;\mathbb Q)=0$, both blocks in the collage of the
graph $B\mathcal C|_{\{0,1\}}$ with $BM$ have zero second
cohomology.  The relative suspension quotient then kills all
degree-two products.  Assume $H^2(BM;\mathbb Q)\ne0$.

The minimal-ideal augmentation description forces at least two
minimal rows and at least two minimal columns.  Its maximal group
is trivial: a nontrivial group would make the ideal contain at
least eight elements, in addition to the identity outside that
ideal.  (If the identity were in the minimal ideal, $M$ would
itself be a group and have zero positive rational cohomology.)
Consequently every homomorphism from $M$ to a group is trivial.
Indeed its restriction to the rectangular minimal band is the
identity element of the group, and multiplication by any band
element then forces the same value on every element of $M$.

Again write $\beta_x:T\to U$ for the equivariant composition
maps.  A constant $\beta_x$ gives an $M$-fixed incoming arrow
and hence zero degree-two products.  Otherwise every $\beta_x$
is injective, with two-element invariant image in the three-element
set $U$.  Two distinct images have a singleton invariant
intersection, giving the same fixed-arrow conclusion.
We may thus suppose all images are one common set $W$.
If two of the maps differ, their ratio is the transposition of
$W$, so the $M$-action on $W$ commutes with that transposition.
Every such actor is a permutation; the preceding triviality of
group quotients makes the action trivial.  This again supplies
a fixed incoming arrow.

The only remaining case has all $\beta_x$ equal.  Collapse $X$
to $x_0$.  The explicit normalized chain-kernel calculation in
Proposition~\ref{C:two-small-incoming-homs} does not require
surjectivity of $\beta_{x_0}$: the two faces still cancel because
all composition maps coincide.  If the action category of $T$
is disconnected, its two states and therefore their images in
$U$ are fixed.  Otherwise the collapse induces an isomorphism
on rational second cohomology.  Its target deletes the middle
object using $t\mapsto tx_0$, leaving the ordinary $U$-act cone.
Every ordinary act cone on a monoid of order at most nine has
all rational degree-two products zero, by
Lemma~\ref{C:nine-unbounded-act-cups}.  Pullback proves the
asserted square vanishing; polarization also gives all mixed
degree-two products.
\end{proof}

\begin{remark}\label{C:two-small-incoming-homs-scope}
The propositions and their opposites remove, in particular, the
seventeen-arrow linear profiles $(1,1,8;3,2,2)$,
$(1,1,8;2,2,3)$, and $(1,1,7;4,2,2)$.
The first two are the only remaining order-eight corner profiles.
This observation is separate from the frozen
154-profile computation: no running source stream is altered by
this note.
\end{remark}
